% Solution: MWG (1995) Exercise 8.D.3 (first-price auction, two bidders, dollar bids). Input by \solitem.
% Check: check_L2.py.
\begin{exsol}{MWG Exercise 8.D.3}{263}{Advanced}{4}
\asked In a first-price sealed-bid auction with two bidders with known valuations $v_1, v_2$ and
bids in whole dollars: (a) is any strategy strictly dominated? (b) is any weakly dominated?
(c) is there a Nash equilibrium, what is it, and is it unique?

\method Dominance (Lecture~1) and Nash equilibrium for a game with discrete strategies and
discontinuous payoffs (Lecture~2, slide~14). The undercutting logic is the one of the Bertrand
price game on a grid (Lecture~2, slides~29--39; exercise 5.12): whoever wins by a small margin has no
reason to bid more than one unit above the loser. Assume that $v_1, v_2$ are positive integers and
that bids are the non-negative integers, with no upper limit.

\solution
\textbf{Payoffs.} For bidder $i$ with bid $b_i$ against $b_j$:
\[
  u_i(b_i, b_j) = \begin{cases} v_i - b_i & b_i > b_j,\\[2pt] \tfrac12 (v_i - b_i) & b_i = b_j,\\[2pt] 0 & b_i < b_j.\end{cases}
\]

\textbf{(a) No strategy is strictly dominated.} Take any bid $b$ and any (pure or mixed) strategy
$\sigma$, and let $N$ exceed $b$ and every bid in the support of $\sigma$. Against a rival bid of $N$
both $b$ and every bid of $\sigma$ lose, so each earns $0$: $\sigma$ does not do \emph{strictly} better
than $b$ against $N$. Hence no bid, not even one above the valuation, is strictly dominated.

\textbf{(b) Weak dominance.}
\begin{itemize}
  \item \emph{Every bid $b \ge v_i$ is weakly dominated by the bid $v_i - 1$} (a legal bid because $v_i \ge 1$).
  Compare by the rival's bid $y$: if $y < v_i - 1$, bidding $v_i - 1$ wins for $1 > 0 \ge v_i - b$; if $y = v_i - 1$ it
  ties for $\tfrac12 > 0 \ge v_i - b$ (the bid $b$ wins at a non-positive surplus); if $y > v_i - 1$ it loses and earns $0$,
  while $b$ earns $v_i - b \le 0$, $\tfrac12(v_i - b) \le 0$ or $0$. So $v_i - 1$ is never worse and is
  strictly better when $y \le v_i - 1$.
  \item \emph{For $v_i \ge 2$ the bid $0$ is weakly dominated by the bid $1$:} against $y = 0$, bid $1$ gives $v_i - 1 \ge \tfrac12 v_i$;
  against $y = 1$ it gives $\tfrac12 (v_i - 1) > 0$; against $y \ge 2$ both give $0$.
\end{itemize}
Bidding one's valuation or more is never sensible: at best it wins for no surplus. Bids in $\{1, \dots, v_i - 1\}$
are not weakly dominated by any other bid (checked by brute force below).

\textbf{(c) Nash equilibria.} Equilibrium profiles split into ties and wins, and each is characterised by
the one-unit deviations.

\emph{Ties, $b_1 = b_2 = b$.} Bidder $i$ earns $\tfrac12(v_i - b)$.
Raising the bid by one wins for $v_i - b - 1$, which is not better iff $v_i - b \le 2$. Lowering loses and gives $0$,
which is not better iff $b \le v_i$. Larger changes are worse. So a tie is an equilibrium iff
\[
  \max\{v_1, v_2\} - 2 \;\le\; b \;\le\; \min\{v_1, v_2\} \quad (b \ge 0).
\]
\emph{Bidder 1 wins, $b_1 > b_2$.} If $b_1 > b_2 + 1$, bidder~1 gains by lowering to $b_2 + 1$
($v_1 - b_2 - 1 > v_1 - b_1$), so $b_1 = b_2 + 1$. Lowering to a tie gives $\tfrac12(v_1 - b_2) = \tfrac12 (v_1 - b_1 + 1)$,
not better than $v_1 - b_1$ iff $b_1 \le v_1 - 1$. The loser earns $0$; by tying at $b_1$ he would
earn $\tfrac12(v_2 - b_1)$, not better iff $b_1 \ge v_2$ (bidding higher is worse still). With
$b_2 = b_1 - 1 \ge 0$ this gives the equilibria
\[
  (b_1, b_2) = (k, k - 1), \qquad \max\{v_2, 1\} \le k \le v_1 - 1 .
\]
\emph{Bidder 2 wins:} symmetrically $(b_1, b_2) = (k - 1, k)$ with $\max\{v_1, 1\} \le k \le v_2 - 1$.

\textbf{Examples.} $v_1 = v_2 = 5$: ties at $b \in \{3,4,5\}$; no wins. $v_1 = 7$, $v_2 = 5$: tie at $(5,5)$ and
wins $(5,4)$, $(6,5)$. $v_1 = 8$, $v_2 = 5$: wins $(5,4)$, $(6,5)$, $(7,6)$ only.

\textbf{Existence and uniqueness.} An equilibrium always exists (for $v_1 \ge v_2$ either $v_1 - v_2 \le 2$ and the tie
range or a win is non-empty, or $v_1 - v_2 \ge 3$ and the win range $[v_2, v_1 - 1]$ is non-empty), but \emph{it is never
unique}: with $d = |v_1 - v_2|$ there are $3$ equilibria if $0 \le d \le 2$ except when $v_1 = v_2 = 1$ (two:
$(0,0)$ and $(1,1)$), and $d$ equilibria if $d \ge 3$. In every equilibrium bids are at most $\max\{v_1, v_2\}$.

\answer{(a) No strategy is strictly dominated (against a very high rival bid all bids lose). (b) Every bid
$b_i \ge v_i$ is weakly dominated by $v_i - 1$ (and the bid $0$ by $1$ when $v_i \ge 2$). (c) There are Nash
equilibria, but never a unique one: ties $b_1 = b_2 = b$ with $\max\{v_1,v_2\} - 2 \le b \le \min\{v_1,v_2\}$, and
profiles where one bidder wins by exactly one dollar, $b_{\text{win}} = b_{\text{lose}} + 1$ with
$\max\{v_{\text{lose}}, 1\} \le b_{\text{win}} \le v_{\text{win}} - 1$. The bidder with the higher valuation wins whenever
$v_1 - v_2 \ge 3$.}

\interp The auction is a Bertrand game in reverse: bidders ``undercut'' the object's price upward
until the surplus is exhausted. With a one-dollar grid the competition stops at a one-dollar margin; the
equilibrium price ends between the loser's value and the winner's value minus one, and ties appear
when valuations are close. Bidding above value is weakly dominated, which is the reason the first-price
auction is not strategy-proof: the \emph{amount} of shading below $v_i$ is what the equilibrium determines
(Lecture~7 will return to this with incomplete information).

\verify \texttt{check\_L2.py} computes all Nash equilibria by brute force (exact fractions, bids
$0, \dots, \max\{v_1,v_2\} + 4$) for every pair $1 \le v_1, v_2 \le 10$ and compares them with the characterisation above and with the counting
rule: they coincide in all $100$ cases, and the number of equilibria is at least two in each. It also
verifies (a) by linear programming against a rival bid above the grid, and (b) by enumerating weak
dominance among pure bids for $v = 1, \dots, 8$ (undominated: $\{1, \dots, v - 1\}$, and $\{0\}$ if $v = 1$). The MWG solutions manual
treats this exercise; the derivation here is complete without it.

\mistakes
\begin{itemize}
  \item Claiming the answer to (a) is ``bids above $v_i$'': against a very high rival bid every bid
  earns zero, so strict dominance fails; they are only weakly dominated.
  \item Reporting the single equilibrium $b_1 = b_2 = v_2$ (the continuous-bid intuition): on a
  dollar grid a tie is an equilibrium only if $|v_1 - v_2| \le 2$, and winning equilibria exist as well.
  \item Treating a tie as winning with probability one: a tie pays $\tfrac12 (v_i - b)$.
\end{itemize}
\end{exsol}
